\chapter{作品设计与实现}
\section{总体架构与信任边界}
当前主线由Python可信Host/网关、真实stdio MCP工具、Python控制面和Vue控制台组成。控制台负责创建任务、查看来源和审批；Host维护当前任务和模型循环；网关在实际工具调用之前核验合同、权限、来源和效果，执行器产生可观察的文件、网络、Git或记忆状态。图\ref{fig:architecture}给出实际组件关系。

\begin{figure}[htbp]\centering
\begin{tikzpicture}[node distance=9mm,box/.style={draw,rounded corners=2pt,align=center,text width=0.88\textwidth,inner sep=7pt,font=\normalsize},arr/.style={-{Stealth[length=2mm]},thick}]
\node[box] (ui) {操作者控制台：任务、可信审批、来源图、实际影响与导出};
\node[box,below=of ui] (host) {可信Host：模型调用、角色/会话、合同/Grant、共享预算};
\node[box,below=of host] (gate) {执行网关：来源用途、真实参数、效果映射、策略与提交复核};
\node[box,below=of gate] (tools) {真实MCP与受控执行器：文件、HTTP、注册程序、Git、Memory};
\node[box,below=of tools] (oracle) {独立观察：实际报告、Receiver、Git对象/ref、SQLite版本与收据};
\draw[arr] (ui)--(host);
\draw[arr] (host)--(gate);
\draw[arr] (gate)--(tools);
\draw[arr] (tools)--(oracle);
\end{tikzpicture}
\caption{当前Host与工具执行链。外部正文和模型提议在执行前接受核验，审批来自可信入口。}\label{fig:architecture}
\end{figure}


授权数据由可信程序提供，资料中的“管理员已批准”或“应扩大权限”文本不能改变授权。当前实现以对象、版本哈希及执行前状态检查绑定动作；签名策略分发、持久化任务和多操作者服务属于后续产品化工作。

\section{任务契约与权限收窄}
TaskContract确定目标、资源范围、允许的目的与副作用；TaskGrant补充主体、有效期、预算和工具权限。对一次动作，网关综合主体权限、合同和当前Grant判定：
\begin{equation}
\mathcal A_{\mathrm{effective}}=\mathcal A_{\mathrm{principal}}\cap\mathcal A_{\mathrm{contract}}\cap\mathcal A_{\mathrm{grant}}.
\end{equation}
子任务在父权限范围内收窄资源、操作和预算，父撤销与角色生命周期进入检查。模型或工具返回不能签发新Grant。预算在任务、角色和祖先scope之间核验，避免每个子角色重新获得一份完整预算。

\section{来源用途与模型请求绑定}
每份Artifact关联实际内容哈希、父来源、可信等级、敏感等级和允许用途。受控复制与已支持的确定性转换保留对应关系；摘要、自由生成和子任务返回采用保守父来源集合。派生内容取父来源中较高的敏感度和较窄的允许用途，普通文本变换不能将资料升级为授权。

对来源父集合$P$，允许用途可概括为：
\begin{equation}
U_{\mathrm{derived}}\subseteq\bigcap_{p\in P}U_p.
\end{equation}
完整原文可用于受限分析和合法引用。公开IPv4 IOC由可信提取器计算固定格式结果，仍按任务范围和审批控制传输。该提取器的能力限定于所登记的公开IOC格式。

模型输入保存messages/tools快照、调用和角色标识；当前主线补充有效HTTP正文与受检输入的绑定证据。字面和JSON转义可以建立确定区间，本地固定分词器可给出可复现token位置。自由生成使用父来源集合；本地映射与云端隐藏模板分别表述。

\section{效果适配、策略与三态决策}
效果适配器将物化后的工具参数转换为SemanticAction，包括资源、目的地、效果类型、输入构件和参数哈希。策略据此检查动作，而不只依赖工具名。当前DSL使用受限表达式和已知事实命名空间，输出ALLOW、ASK或BLOCK；硬拒绝优先于普通拒绝、审批与放行，语义检测不能降低确定性拒绝等级。

文件写入、封闭HTTP请求、注册程序执行、Git发布、记忆和委派分别由适配器限制实际参数。注册程序以程序身份、参数和摘要核验，并非任意Shell执行。Git检视保持远端不变；授权发布绑定确切commit、remote、ref与预期旧值。网络外发由封闭Receiver给出实际接收记录。

\section{审批与执行前复核}
IOC与Git业务可在ASK处暂停，操作者先核对动作预览，再批准和提交。批准本身不执行动作。提交时重新核对实际输入、来源、工具版本、合同、Grant、策略、资源状态和预算；相关对象漂移后旧批准失效。确认丢失时记录UNKNOWN并停止自动重发。

\begin{figure}[htbp]\centering
\begin{tikzpicture}[node distance=8mm and 5mm,box/.style={draw,rounded corners=2pt,align=center,text width=0.25\textwidth,inner sep=6pt,font=\normalsize},wide/.style={box,text width=0.86\textwidth},arr/.style={-{Stealth[length=2mm]},thick}]
\node[wide] (proposal) {模型提议：物化参数，绑定任务、来源与工具版本};
\node[wide,below=of proposal] (policy) {确定性授权、来源用途、DSL与风险证据};
\node[box,below=of policy] (ask) {ASK\\可信审批后提交};
\node[box,left=of ask] (block) {BLOCK\\记录原因与安全反馈};
\node[box,right=of ask] (allow) {ALLOW\\进入执行前复核};
\node[wide,below=of ask] (commit) {提交复核：输入、版本、权限、预算与资源状态仍一致};
\node[wide,below=of commit] (result) {执行并观察实际效果；记录成功、失败或UNKNOWN};
\draw[arr] (proposal)--(policy);
\draw[arr] (policy)--(ask);
\draw[arr] (policy.west)-|(block.north);
\draw[arr] (policy.east)-|(allow.north);
\draw[arr] (ask)--(commit);
\draw[arr] (allow.south)|-(commit.east);
\draw[arr] (commit)--(result);
\end{tikzpicture}
\caption{三态决策、动作绑定审批与执行前复核。BLOCK分支不进入执行。}\label{fig:decision}
\end{figure}


图\ref{fig:decision}区分审批、提交和实际结果。进程内单次提交约束防止同一批准被重复消费；跨服务重启的持久审批与一致性尚需独立实现。

\section{多层Agent、记忆与业务保全}
委派业务运行main、review、verify三个实际对话主体，子角色以收窄的Grant读取指定资料，父Host接收保留来源的结果。当前角色在同一Host内，操作系统层面的角色隔离另行验收。

记忆更新包含capture与reconcile两个实际会话。旧原文保留为revision 1，新会话读取旧记忆与当前资料，经CAS追加revision 2并复读；独立Oracle检查版本、原文、previous hash与主体/会话顺序。记忆正文只提供证据，授权记录由可信程序维护。

阻断反馈包含错误、受限效果和可继续执行的信息。报告入口还可使用无工具分析角色，Host核对引用位置和原文字面覆盖。能否继续完成任务由实际报告事实决定；读取失败或重复无效提议按有界Host规则结束。通用恢复策略的相对收益仍需比较实验。

\section{部署与当前运行方式}
开发环境需要Python 3.12、前端Node 20.19以上或22.12以上，以及Git业务所需的Git。依赖由lock文件固定。无密钥可运行确定性演示、工具与离线核验；云模型由可信Provider加载仓库外凭据，并受共享请求预算约束。

M23/M24另有Docker Desktop Linux容器smoke。M24六条链在非root、只读rootfs、无外网、移除capabilities和资源限额下运行；同一容器和UID承载Host及工具。它验证指定拓扑的有限运行性，完整多服务部署、独立角色隔离和性能验收另行安排。现行控制面为Python，Java/MySQL不是运行依赖。
